% ============================================================
%  ch9-end.tex — Implementation, Conclusion, References
% ============================================================

\chapter{System Implementation}
\label{ch:impl}

\section{UAV Platform}
An Unmanned Aerial Vehicle (UAV), commonly a \emph{drone}, is an aircraft without
a human pilot on board and is one component of an Unmanned Aircraft System (UAS),
which also includes a ground controller and a communications link. Flight may be
autonomous or remotely piloted. A quadcopter generates lift from four rotors
working together, each lifting roughly a quarter of the weight, which allows
smaller, cheaper motors; movement is controlled by varying the power each motor
delivers.

\begin{table}[H]
\centering
\caption{Airframe components of the quadcopter platform.}
\label{tab:parts}
\rowcolors{2}{mist}{white}
\begin{tabularx}{\textwidth}{@{}l X@{}}
\toprule
\rowcolor{navy}
\textcolor{white}{\textbf{Component}} & \textcolor{white}{\textbf{Purpose}}\\
\midrule
Frame & Centre plate, four arms and four motor brackets; aluminium / carbon fibre for strength with vibration tolerance. \\
Motors & Four brushless motors; lighter on the battery. \\
ESCs & Electronic speed controllers deliver power to the motors; one per arm. \\
Propellers & Matched to the frame; plastic / carbon fibre, correctly fitted. \\
Battery + monitor & Chosen by capacity and type; monitor warns of low charge. \\
Autopilot & Pixhawk-class flight controller running PX4 / ArduPilot. \\
Companion computer & Single-board computer hosting the sensing payload. \\
RC link & Receiver for transmitter commands and telemetry. \\
GPS + compass & Positioning for navigation and for geo-tagging samples. \\
Camera + gimbal & Optional stabilised capture for aerial imagery. \\
\bottomrule
\end{tabularx}
\end{table}

The flight stack (PX4 / ArduPilot on a Pixhawk-class autopilot, with a companion
computer for the payload) flies the planned coverage tour of \Cref{ch:algo} as a
waypoint mission. Operating a sensing payload over a campus is subject to
airspace rules: the aircraft broadcasts \textbf{Remote ID}, stays within
authorised, geofenced airspace, and is registered where required (in India, under
the DGCA Drone Rules 2021 and the Digital Sky platform). These constraints are
part of the mission design rather than an afterthought.

\section{Sensing Payload}
The payload is a single-board computer (Raspberry Pi 5 / 4) carrying an external
monitor-mode wireless adapter and a GPS module. The on-board Broadcom radio has
limited monitor-mode support, so an external adapter with a monitor-capable
chipset is used --- a MediaTek MT7612U for dual-band 802.11ac, or a Wi-Fi~6/6E
adapter for 6\,GHz capture. The GPS supplies the position that tags every captured
frame so that samples can be placed on the survey grid.

\subsection{Access-point / hotspot configuration}
For field control the companion computer exposes its own management hotspot,
sharing a wired or cellular uplink over its wireless card. On current Raspberry Pi
OS this is configured through \textbf{RaspAP} or the \code{nmcli device wifi
hotspot} command (NetworkManager), which assigns the interface address, runs
DHCP/DNS, and enables IP forwarding and NAT between the uplink and the hotspot in
a single managed step.

\subsection{Monitor-mode capture}
A monitor interface is created on the external adapter and the capture is written
to a PCAP for analysis, or parsed live with \code{Scapy}:
\begin{lstlisting}[style=shell]
# Create a monitor interface on the external adapter
sudo airmon-ng start wlan1

# Option A: capture to a file for later analysis
sudo tcpdump -ni wlan1mon -s 0 -w survey.pcap

# Option B: parse beacons live in Python (passive; management frames only)
sudo python -m uav_sniffer.cli analyse --pcap survey.pcap \
     --box 17.470 78.360 17.474 78.365
\end{lstlisting}
The captured 802.11 beacon frames (source BSSID, SSID, channel, advertised
security and RSSI) are exactly the records required by the dead-zone and rogue-AP
logic of \Cref{ch:dead,ch:rogue}.

\section{System Specification Summary}
\Cref{tab:spec} summarises the components that make up the system.

\begin{table}[H]
\centering
\caption{System specification summary.}
\label{tab:spec}
\footnotesize
\rowcolors{2}{mist}{white}
\begin{tabularx}{\textwidth}{@{}>{\raggedright\arraybackslash}p{0.26\textwidth} >{\raggedright\arraybackslash}X@{}}
\toprule
\rowcolor{navy}
\textcolor{white}{\textbf{Subsystem}} & \textcolor{white}{\textbf{Choice}}\\
\midrule
Flight platform     & Quadcopter, Pixhawk-class autopilot, PX4 / ArduPilot \\
Companion computer  & Raspberry Pi 5 / 4 (Raspberry Pi OS, NetworkManager) \\
Capture radio       & External monitor-mode adapter (MT7612U / Wi-Fi 6E) \\
Spectrum            & 2.4, 5 and 6\,GHz (Wi-Fi 6E); extensible to Wi-Fi 7 \\
Capture stack       & Scapy / Kismet / \code{tshark}; PCAPNG + JSON \\
Path planning       & Boustrophedon-seeded SOM with point-in-polygon culling \\
Dead-zone model     & RSSI/SNR heatmap (IDW/RBF) with dead \& weak contours \\
Rogue detection     & Allow-list + evil-twin + security-capability + hidden-node \\
Analytics           & MQTT stream $\rightarrow$ time-series $\rightarrow$ live dashboard; batch retention \\
Positioning         & On-board GPS geo-tagging per frame \\
Compliance          & Remote ID, DGCA Drone Rules 2021 / Digital Sky, geofencing \\
Privacy             & Salted-hash device identifiers; least-data retention \\
\bottomrule
\end{tabularx}
\end{table}

\section{Security, Privacy and Ethics}
\label{sec:ethics}
Because the system captures wireless traffic and device identifiers, it is
operated under a strict authorised-use policy:
\begin{itemize}
  \item Capture only on networks you own or are contractually authorised to test;
        keep written scope and authorisation.
  \item Survey \emph{passively} only --- read broadcast management frames; never
        associate, inject, deauthenticate, or decode third-party payloads.
  \item Treat MAC/BSSID addresses as personal data --- hash or truncate them,
        minimise retention, and account for client MAC randomization when counting
        devices.
  \item Comply with local radio, aviation and privacy law (DGCA Drone Rules,
        Remote ID, and applicable data-protection obligations).
\end{itemize}
A man-in-the-middle interception is powerful evidence of the \emph{risk} a rogue
AP poses, but it is never performed against third-party traffic; the accompanying
software ships no credential-capture functionality and defaults to passive
surveying.

% ============================================================
\chapter{Conclusion}
This report presented an end-to-end system in which a UAV performs complete-area
coverage of a region, passively sniffs 802.11 traffic along an obstacle-aware path
planned by a boustrophedon-seeded Self-Organising Map, and uses the resulting
signal map to detect wireless dead-zones and rogue access points, supported by
alert/feedback and real-time analytics subsystems. The coverage experiments
showed that path length scales predictably with node count and that a single
obstacle changes run-time only marginally while leaving path length essentially
unchanged --- confirming that obstacle-aware culling integrates cleanly with the
planner.

The detection pipeline treats coverage as a measured signal surface rather than a
binary test, surfacing weak zones before they fail, and classifies access points
by fusing an allow-list with evil-twin, security-capability and hidden-node
evidence so that detection survives WPA3 protected management frames and client
MAC randomization. The whole pipeline --- planning, passive capture across the
2.4/5/6\,GHz bands, heatmap dead-zone extraction, rogue-AP classification and a
streaming analytics path --- is released as a reproducible, documented
open-source package, operated under a passive, privacy-preserving and strictly
authorised policy.

% ============================================================
\begin{thebibliography}{99}
\addcontentsline{toc}{chapter}{References}
\small
\bibitem{landrobots} \emph{Category: Land robots}, Wikipedia. \url{https://en.wikipedia.org/wiki/Category:Land_robots}
\bibitem{zelinsky_phd} A.~Zelinsky, ``Environment Exploration and Path Planning Algorithms for a Mobile Robot using Sonar,'' Ph.D.\ dissertation, Univ.\ of Wollongong, Dept.\ of Computer Science, Oct.\ 1991.
\bibitem{zelinsky1993} A.~Zelinsky, R.~A.~Jarvis, J.~C.~Byrne, S.~Yuta, ``Planning paths of complete coverage of an unstructured environment by a mobile robot,'' in \emph{Proc.\ Int.\ Conf.\ on Advanced Robotics}, vol.~13, Nov.\ 1993, pp.~533--538.
\bibitem{yun2010} S.~C.~Yun, V.~Ganapathy, L.~O.~Chong, ``Improved genetic algorithms based optimum path planning for mobile robot,'' in \emph{Int.\ Conf.\ on Control, Automation, Robotics \& Vision (ICARCV)}, Dec.\ 2010, pp.~1565--1570, IEEE.
\bibitem{jimenez2007} P.~A.~Jimenez, B.~Shirinzadeh, A.~Nicholson, G.~Alici, ``Optimal area covering using genetic algorithms,'' in \emph{IEEE/ASME Int.\ Conf.\ on Advanced Intelligent Mechatronics}, Sep.\ 2007, pp.~1--5.
\bibitem{tsp_compare} \emph{Comparison of TSP Algorithms}, SlideShare. \url{https://www.slideshare.net/kaalnath/comparison-of-tsp-algorithms}
\bibitem{youtube_tsp} \emph{Travelling Salesman / path heuristics}, video. \url{https://www.youtube.com/watch?v=q6fPk0--eHY}
\bibitem{makeuseof} \emph{What is a wireless dead zone and how to eliminate it}, MakeUseOf. \url{https://www.makeuseof.com/tag/what-is-a-wireless-dead-zone-and-how-to-eliminate-it/}
\bibitem{zone13} \emph{Wi-Fi monitoring using Raspberry Pi}, Zone13. \url{https://zone13.io/post/wifi-monitoring-using-raspberry-pi/}
\bibitem{scapy} P.~Biondi \emph{et al.}, \emph{Scapy: Packet crafting for Python}. \url{https://scapy.net}
\bibitem{kismet} \emph{Kismet: wireless network detector, sniffer and WIDS}. \url{https://www.kismetwireless.net}
\bibitem{dgca} Ministry of Civil Aviation (India), \emph{Drone Rules 2021} \& Digital Sky platform.
\bibitem{faa_rid} U.S.\ FAA, \emph{Remote Identification of Unmanned Aircraft}, Final Rule, 14 CFR Part 89.
\end{thebibliography}
